%%%%%%%%%%%%%%%%%%%%%%%%%%%%%%%%%%%%%%%%%%%%%%%%%%%%%%%
% NOTAS TEORIA DE TRANSPORTE OPTIMO
%
% Contains ?? figures. 
% File started: ??
% First circulated version: ??
% Last version:
% Last modification: -
%%%%%%%%%%%%%%%%%%%%%%%%%%%%%%%%%%%%%%%%%%%%%%%%%%%%%%%

\documentclass[11pt,a4paper]{amsart}
\usepackage{xr-hyper}
\newcounter{chapter}

\usepackage{amssymb}
\usepackage{hyperref}
\usepackage{listings}
\usepackage{mathtools}
\usepackage{enumitem}

% Fuzz -------------------------------------------------------------------
\hfuzz2pt % Don't bother to report over-full boxes if over-edge is < 2pt
% Line spacing -----------------------------------------------------------

%%%%%%%%%fin pegada%%%%%%%%%%%%%%%%%%5

\usepackage[T1]{fontenc}
\usepackage[spanish]{babel}
\usepackage{graphicx}
\usepackage{tikz}
\usetikzlibrary{babel}

\setlength{\textwidth}{15.5cm} 
\setlength{\textheight}{24cm}
\setlength{\oddsidemargin}{0cm}
\setlength{\evensidemargin}{0cm}


\parskip 3pt

%%
%% OPERADORES
%%
\renewcommand{\Re}{\operatorname{Re}}
\newcommand{\tr}{\operatorname{tr}}
\newcommand{\id}{\operatorname{id}}
\def\div{\operatorname {\mathrm{div}}}
\def\Dom{\operatorname {\mathrm{Dom}}}
\def\argmin{\operatorname {\mathrm{arg\, min}}}
\def\argmax{\operatorname {\mathrm{arg\, max}}}
\def\gen{\operatorname {\mathrm{gen}}}
\def\Im{\operatorname {\mathrm{Im}}}
\def\Nu{\operatorname {\mathrm{Nu}}}
\def\sop{\operatorname {\mathrm{sop}}}
\def\diam{\operatorname {\mathrm{diam}}}
\def\dist{\operatorname {\mathrm{dist}}}
\def\var{\operatorname {\mathrm{Var}}}
\def\codim{\operatorname {\mathrm{codim}}}
\def\pint{\operatorname {--\!\!\!\!\!\int\!\!\!\!\!--}}

%%
%% EPSILON
%%
\newcommand{\ve}{\varepsilon}

%%
%% LETRAS MATHBB
%%
\newcommand{\R}{{\mathbb R}}
\newcommand{\Z}{{\mathbb Z}}
\newcommand{\C}{{\mathbb C}}
\newcommand{\Q}{{\mathbb Q}}
\newcommand{\N}{{\mathbb N}}
\newcommand{\E}{{\mathbb E}}
\newcommand{\V}{{\mathbb V}}
\renewcommand{\P}{{\mathbb P}}
\newcommand{\II}{{\mathbb I}}
\newcommand{\OO}{{\mathbb O}}
\newcommand{\Sb}{{\mathbb S}}
\newcommand{\Tb}{{\mathbb T}}

%%
%% LETRAS CALIFGRAFICAS
%%
\newcommand{\F}{{\mathcal F}}
\newcommand{\Hc}{{\mathcal H}}
\newcommand{\Tc}{{\mathcal T}}
\newcommand{\Ec}{{\mathcal E}}
\newcommand{\Vc}{{\mathcal V}}
\newcommand{\A}{{\mathcal A}}
\newcommand{\J}{{\mathcal J}}
\newcommand{\G}{{\mathcal G}}
\newcommand{\U}{{\mathcal U}}
\newcommand{\B}{{\mathcal B}}
\newcommand{\Sw}{{\mathcal S}}
\newcommand{\D}{{\mathcal D}}
\newcommand{\M}{{\mathcal M}}
\newcommand{\LL}{{\mathcal L}}

%%
%% LETRAS MATHBF
%%
\newcommand{\ab}{{\mathbf a}}
\newcommand{\bb}{{\mathbf b}}
\newcommand{\cc}{{\mathbf c}}
\newcommand{\ee}{{\mathbf e}}
\newcommand{\ff}{{\mathbf f}}
\newcommand{\gb}{{\mathbf g}}
\newcommand{\hh}{{\mathbf h}}
\newcommand{\n}{{\mathbf n}}
\newcommand{\mm}{{\mathbf m}}
\newcommand{\pp}{{\mathbf p}}
\newcommand{\qq}{{\mathbf q}}
\newcommand{\sbf}{{\mathbf s}}
\newcommand{\tb}{{\mathbf t}}
\newcommand{\uu}{{\mathbf u}}
\newcommand{\vv}{{\mathbf v}}
\newcommand{\ww}{{\mathbf w}}
\newcommand{\xx}{{\mathbf x}}
\newcommand{\yy}{{\mathbf y}}
\newcommand{\zz}{{\mathbf z}}
\newcommand{\I}{{\mathbf 1}}
\newcommand{\Pb}{{\mathbf P}}
\newcommand{\XX}{{\mathbf X}}
\newcommand{\YY}{{\mathbf Y}}

%%
%% LETRAS MATHFRAK
%%

\newcommand{\Gk}{{\mathfrak G}}


%%
%% TEOREMAS
%%
\newtheorem{teo}{Teorema}[section]
\newtheorem{lema}[teo]{Lema}
\newtheorem{prop}[teo]{Proposición}
\newtheorem{cor}[teo]{Corolario}

%%
%% REMARK
%%
\theoremstyle{remark}
\newtheorem{obs}[teo]{Observación}

%%
%% DEFINICION
%%
\theoremstyle{definition}
\newtheorem{defi}[teo]{Definición}
\newtheorem{ejem}[teo]{Ejemplo}
\newtheorem{ejer}[teo]{Ejercicio}
\newtheorem{nota}[teo]{Notación}
\newtheorem{ejercicio}{Ejercicio}[chapter]


%%
%% NUMERACION
%%
\numberwithin{subsection}{section}
\numberwithin{equation}{section}

%%
%% TITULO
%%


\externaldocument[N-]{../notas/Notas-TO}
\newcommand{\cuadernolink}[2]{\noindent\textbf{Cuaderno:} \emph{#1} (\href{https://colab.research.google.com/github/jfbonder/transporte-optimo/blob/main/cuadernos/#2}{abrir en Colab}).\par\medskip}
\pagestyle{plain}

\begin{document}
\renewcommand{\thechapter}{B}
\begin{center}
{\large\sc Teoría de Transporte Óptimo}\\[2pt]
{\small 2do cuatrimestre 2026 --- Julián Fernández Bonder}\\[10pt]
{\LARGE\bf Ejercicios del Apéndice B}\\[4pt]
{\large El teorema de Fenchel--Rockafellar}
\end{center}
\bigskip
\noindent{\small Los ejercicios son los de la sección de ejercicios del Apéndice~B de las notas del curso, con la misma numeración. Las referencias a teoremas, proposiciones y secciones remiten a las notas (\url{https://github.com/jfbonder/transporte-optimo}).}
\medskip


\begin{ejercicio}[Supremos de funciones convexas y semicontinuas]\label{ejer.FR.sup}
Sea $E$ un espacio normado y $\{f_\alpha\}_{\alpha\in A}$ una familia de
funciones $E\to\R\cup\{+\infty\}$. Probar que si todas las $f_\alpha$ son
convexas, $f=\sup_\alpha f_\alpha$ es convexa, y que si todas son
semicontinuas inferiormente, $f$ lo es. Deducir que la transformada de
Legendre--Fenchel $F^*$ de cualquier función $F$ es convexa y semicontinua
inferiormente (para la topología débil-$*$ de $E^*$, y en particular para la
de la norma).
\end{ejercicio}

\begin{ejercicio}[Ejemplos de transformadas]\label{ejer.FR.ejemplos}
Calcular la transformada de Legendre--Fenchel de las siguientes funciones
de $\R$ en $\R\cup\{+\infty\}$:
\begin{enumerate}[label=(\alph*)]
\item $f_p(x)=\frac{x^p}{p}$ para $x\ge0$ y $+\infty$ para $x<0$, con $p>1$
(se obtiene $f_{p'}$ con $\frac1p+\frac1{p'}=1$, restringida a $y\ge0$);
\item $f_\infty(x)=0$ si $x\in[0,1]$, $+\infty$ en otro caso; interpretar el
resultado como límite $p\to\infty$ de (a);
\item $f(x)=|x|$, $f(x)=e^x$, $f(x)=-\log x$ ($x>0$), $f(x)=x\log x$ ($x\ge0$).
\end{enumerate}
Verificar en cada caso que $f^{**}=f$.
\end{ejercicio}

\begin{ejercicio}[Mínimos locales y globales]\label{ejer.FR.minimo}
Sea $C\subset E$ convexo y $f\colon C\to\R\cup\{+\infty\}$ convexa y propia.
Probar que todo mínimo local de $f$ es global, y que el conjunto de
minimizadores es convexo. Dar un ejemplo de una función convexa en $\R^2$
con infinitos minimizadores, y otro sin minimizadores pero con ínfimo finito.
\end{ejercicio}

\begin{ejercicio}[Biconjugada]\label{ejer.FR.biconjugada}
Sea $F\colon E\to\R\cup\{+\infty\}$ propia.
\begin{enumerate}[label=(\alph*)]
\item Probar que $F^{**}\le F$ y que $F^{**}$ es la mayor función convexa
y semicontinua inferiormente que está por debajo de $F$.
\item Probar que si $F$ es convexa, semicontinua inferiormente y propia,
entonces $F^{**}=F$. \emph{Sugerencia:} si $F^{**}(x_0)<F(x_0)$, separar el
punto $(x_0,F^{**}(x_0))$ del epigrafo de $F$ (cerrado y convexo en
$E\times\R$) mediante Hahn--Banach, como en la demostración del
Teorema~\ref{N-teo.FR}.
\item Calcular $F^{**}$ para $F(x)=\min\{|x-1|,|x+1|\}$ en $\R$.
\end{enumerate}
\end{ejercicio}

\begin{ejercicio}[El caso de dimensión finita]\label{ejer.FR.dimension.finita}
Sea $K\subset\R^n$ convexo cerrado no vacío y $\chi_K$ su función indicatriz
($0$ en $K$, $+\infty$ fuera).
\begin{enumerate}[label=(\alph*)]
\item Probar que $\chi_K^*(y)=\sup_{x\in K}x\cdot y$ es la \emph{función
soporte} de $K$, y que $K=\{x:\ x\cdot y\le\chi_K^*(y)\ \forall y\}$.
\item Aplicar el Teorema~\ref{N-teo.FR} con $F=\chi_K$ y $G(x)=c\cdot x$
para reobtener la dualidad de programación lineal del
Capítulo~\ref{N-cap.PL} (con $K=\{x\ge0:\ Ax=b\}$), identificando qué
hipótesis del teorema corresponde a la factibilidad del primal.
\end{enumerate}
\end{ejercicio}

\end{document}
